% !TeX root = ../papers/nc-p3.tex
\section{中心商を固定した三重み全体の正則性}
\label{sec:half-period-weight-fibers}
前節の枠で $K=\langle p_1,p_2,p_3,p_4\rangle$ と置く。
重みは $\omega_i\ne0$, $\sum\omega_i=1$ とし、$I\cap L_+=I\cap L_-=0$ である範囲を取る。
これは全ての隣接階数が6となる条件である。一文字切断は全階数である。
実際各 $L_\pm$ が可逆な交代形式を含むため、各辺の二方向の台はともに全四空間となる。
折り畳んだ二頂点代数を $A_\omega$、中心二次切断を $q$ と書く。
\begin{theorem}\label{thm:half-period-weight-fibers}
共通の中心商 $B$ は関係 $K+L_+$、$(K+L_-)^{\mathrm{op}}$ を持つ。
次の二条件は同値である。
\begin{enumerate}
\item $B$ は各始点で $H_B=(1-t^2)/(1-t)^4$ を持ち、左右の次数付きsocleが零である。
\item 全ての上記重みで $A_\omega$ は標準Koszul左右AS正則で $q$ は非零因子である。
\end{enumerate}
重みの範囲が空でないとする。この範囲のただ一つの重みで後者の正則性と非零因子性が
成立すれば両条件が従う。従ってこの判定には枠の六自由度だけが残る。
\end{theorem}
\begin{proof}
前節の積表示から $q$ は中心であり、$I+kq=K$ なので商は重みに依存しない。
全枠・重みの表示空間は既約である。実際、順序付き直線枠と基底空間の可逆変換は既約であり、
滑らかな形式と直線族の順序の対も固定した一対の $\mathrm{GL}_4$ 軌道である。
重みの超平面と上記の階数開条件を加えても既約である。
この族には正則点がある。定理\ref{thm:half-period-frame}の53の窓の枠と等重みの
特性零持上げが一例であり、下記の有理例で直接確かめる。
各内部次数における双対bar複体の議論を適用すれば、
この表示空間の幾何学的一般点で標準Koszul性と標準Hilbert関数が成立する
（命題\ref{prop:unrestricted-good-reduction}の縮約下界と同じ議論を局所の小行列上で行う）。
元の各経路次元の上半連続性により、全特殊点で標準Hilbert関数は下界となる。

第一条件の下で、$q$ による乗法と商への全射から、係数ごとに
$H_{A_\omega}\le H_B/(1-t^2)=(1-t)^{-4}$ である。下界と一致し、
乗法の核が全次数で零である。左右に同じ議論が使える。
中央の非零因子 $q$ の長さ1の自由解消により
$\operatorname{Ext}_{A_\omega}^1(S,A_\omega)
\simeq\operatorname{Hom}_B(S,B)(2)=0$、また $\operatorname{Hom}_{A_\omega}(S,A_\omega)=0$ となる。
ここで $S$ は各頂点単純である。前者は $A_\omega\to B$ の環変更スペクトル列で
$\operatorname{Ext}_{A_\omega}^j(B,A_\omega)$ が $j=1$ の $B(2)$ だけとなることから従う。
左右で同じである。

四次テンソルの縮約を項に持つ長さ4の線型列を取る。
両隣接階数6と一文字階数4から中間係数空間の次元は $1,4,6,4,1$ で、
反対側の表示の最初の二微分を双対化したものが最後の二微分である。
実際、相補縮約間の対合は同じテンソルの切断の左右像の対合であり、
三次から二次への含有は左の二次から一文字への含有に転置する。
関係による表示で最初の二項は完全、上の $\operatorname{Hom}$ と $\operatorname{Ext}^1$ の
消滅で最後の二項も完全である。残る中央では $H_{A_\omega}=(1-t)^{-4}$ の
Euler等式によりホモロジーが零となる。左右に適用して標準Koszul左右AS正則性を得る。
逆に一つの正則重みで $q$ が非零因子なら、乗法の短完全列が所定の商の級数を与え、
同じ環変更同型とAS消滅で商の両socleが零となる。
\end{proof}

\begin{corollary}\label{cor:half-period-rational-frame}
次の有理枠に対し、代数閉体の標数零で非零の全重み $\sum\omega_i=1$ が
標準Koszul左右AS正則な窓を与える。$q$ の対角と各 $u_i$ の行は
\[
 (37,15,47,9/(37\cdot15\cdot47)),\qquad
 \begin{pmatrix}
 1&-5&25&-26\\1&-5&-25&26\\1&5&25&26\\1&5&-25&-26
 \end{pmatrix}.
\]
直線族の順序は標準体積に対して $\sqrt{\det q}=3$ の符号で取る。
同じ剰余枠・直線族順序を持つ任意の非退化な53進持上げでも、
その幾何学的分数体の全ての非零重みに対して成立する。
\end{corollary}
\begin{proof}
$\det q=9$ で $K$ から対角二次式への制限は剰余で階数4である。
従って $K\cap \bigwedge^2E^*=0$ で全重みが二つの階数開条件を満たす。
\path{computations/verify_half_period_weight_fibers.py} は等重みで $h=45$, $c=15$ とした
第一テンソルの53での還元が、前節で基底を揃えたWalsh窓と一致することを直接確かめる。
命題\ref{prop:unrestricted-good-reduction}により等重みの持上げは全次数で正則である。

剰余で $q$ の非零因子性も確認する。Walshの四正規二次式のうち添字3のものは
交換符号 $\phi_3(i)=(-1)^{i_0i_1}=J(i)$ を持つため、この位数2持上げと交換する。
その平方は四中心変数の一つなので、左右とも非零因子である。
各位置で動く十二点の二次評価の階数は9で、両側とも十次元経路商からの核は直線である。
この正規式も、構成した $q$ もこの核に入り、いずれも非零である
（$q$ は固定八点で非零）。従って $q$ は各位置でこの正規式の非零倍となる。
等重みの整持上げの各経路加群は自由であり、剰余で単射な乗法の小行列を持ち上げれば
一般側でも全次数で $q$ は左右非零因子である。定理\ref{thm:half-period-weight-fibers}により
重みへの付値条件は消え、分母または分子が53で消える重みも含む。
\end{proof}
